\chapter{Methodology}\label{chapter:methodology}
This chapter presents the ALERT framework developed in this thesis. ALERT addresses the two main failures of online trackers: identity switches and track fragmentation. ALERT consists of two stages: a tracklet splitter that detects identity switches within tracklets and a tracklet connector that reconnects tracklets with the same identity. Both stages utilize domain-specific attributes such as jersey numbers and team identity, which are predicted after the tracking stage, and geometric and temporal spatial information. Three connector frameworks have been developed and tested, such as a rule-based connector and two learning-based connectors using gradient boosted trees and a transformer.

\

Figure \ref{fig:overall_pipeline} provides an overview of the entire pipeline. As input to the pipeline, we use the newly constructed SoccerNet Player-Tracking dataset. Initial tracklets are generated by running Deep-EIoU on the ground truth detections extracted from the dataset annotations. Jersey numbers and team identity are then predicted for each resulting tracklet. These domain-specific attributes in combination with appearance features and geometrical properties are then used by the tracklet splitter to detect identity switches within tracklets and split those tracklets into multiple pure tracklets. Finally, a tracklet connector reconnects fragmented tracklets belonging to the same player into one final tracklet. Three tracklet connectors have been developed and evaluated: a rule-based connector, a gradient boosted trees connector, and a transformer-based connector.

\

% The chapter is organized as follows: Section \ref{dataset_construction} describes the dataset construction procedure. Section \ref{detection_and_tracking} describes the detection and tracking stage. Section \ref{reid_spatio_temporal_splitter} describes the spatio-temporal ReID splitter. Section \ref{attribute_prediction} describes the attribute prediction pipeline. Section \ref{tracklet_splitter} describes the attribute-based tracklet splitter, and Section \ref{tracklet_merger} presents the three merger designs.

\begin{figure}[h!]
\centering
\includegraphics[width=1\textwidth]{figures/alert_overview.drawio.png}
\caption{Overview of the ALERT pipeline}
\label{fig:overall_pipeline}
\end{figure}

\section{Dataset construction}\label{dataset_construction}
The framework is evaluated on the newly constructed SoccerNet Player-Tracking dataset. This dataset is derived from the SoccerNet Game State Reconstruction dataset (GSR), which contains 164 sequences of 30 seconds of broadcast football footage recorded at 25 frames per second with a resolution of 1920$\times$1080. The dataset is divided into train, validation, and test splits containing 57, 58, 49 sequences respectively. Each sequence is annotated with bounding boxes and track identities for all players, goalkeepers, referees, and the ball. 

\

The post-processing framework utilizes jersey numbers and team identity as signals for splitting and merging. Since referees, the pitch and the ball do not have these attributes, they are excluded from the dataset. The resulting SoccerNet Player-Tracking dataset contains only player and goalkeeper tracklets, and is obtained by filtering the GSR annotations and reformatting them into the MOT-Challenge convention used by the SoccerNet tracking benchmarks.

\paragraph{Source Format}
Each sequence of the SoccerNet Game State Reconstruction dataset is in a separate folder, which contains a directory of frame images and a single annotation file \texttt{Labels-GameState.json}. This JSON file contains all the annotations, where the relevant fields are:
\begin{itemize}
    \item \textbf{Images}. A list  of image entries, each containing an \texttt{image\_id} and a \texttt{file\_name} in the form \texttt{000001.jpg}, where the filename corresponds to the frame number, starting from 1.

    \item \textbf{Annotations}. A list of detection entries, each containing an \texttt{image\_id}, \texttt{track\_id}, \texttt{category\_id}, \texttt{attributes} containing the role, jersey number and team identity, and a bounding box \texttt{bbox\_image} given as $\{x,y,w,h\}$, where $x,y$ is the pixel coordinate of the top-left corner and $w,h$ are the width and height.

    \item \textbf{Categories}. A list that maps the category identifiers to labels. Category 1 corresponds to player, category 2 to goalkeeper, and referees, staff and the ball are assigned to higher identifiers. 
\end{itemize}

\paragraph{Target Format}
The SoccerNet GSR format needs to be converted to the standard MOT format for evaluation. This format contains a single \texttt{.txt} file where each line represents one annotation:
\begin{equation*}
    \texttt{frame}, \texttt{id}, \texttt{x}, \texttt{y}, \texttt{w}, \texttt{h}, \texttt{conf}, -1, -1, -1
\end{equation*}

The \texttt{frame} field is the frame number, \texttt{id} is the ground truth track identity, $x,y$ are the top-left corner of the bounding box, and $w,y$ is its width and height. The \texttt{conf} field stores the detection confidence, which is always 1 since the annotations are the ground truth. The three last fields correspond to 3D position coordinates that are not used and set to -1. Each sequence additionally contains a \texttt{seqinfo.ini} file, denoting the sequence name, frame rate, sequence length, and image dimensions, required by the MOT-Challenge evaluation toolkit.   

\paragraph{Dataset Conversion}
In order to convert the SoccerNet GSR dataset into our newly created SoccerNet Player-Tracking dataset, the following steps need to be applied:

\begin{enumerate}
    \item Parse \texttt{Labels-GameState.json} and build a lookup table that maps each  \texttt{image\_id} to its frame number by taking the filename stem.

    \item Iterate over all annotations  and retain only entries with a \texttt{category\_id} equal to player or goalkeeper. 

    \item Convert each remaining annotation into a MOT row using the lookup table from step 1 and the bounding-box fields from \textbf{bbox\_image}, setting \texttt{conf} to 1 and the three trailing fields to -1.

    \item Sort the resulting  rows by frame and track identity, and write them to a \texttt{gt.txt} file, together with a \texttt{seqinfo.ini} file.

    \item Rename the sequence folder from \texttt{SNGS-XXX} to \texttt{SNPT-XXX} to indicate that the sequence is part of the SoccerNet Player-Tracking dataset.
\end{enumerate}

Applying this procedure to all sequences in the SoccerNet GSR dataset yields the final SoccerNet Player-Tracking dataset used throughout this thesis.

\section{Detection and Tracking}\label{detection_and_tracking}
Detecting players and tracking them across frames is the first part of the pipeline. Instead of using a detection model like YOLO \cite{yolo} or RF-DETR \cite{rfdetr}, we use the ground truth bounding box annotations from the SoccerNet Player-Tracking dataset. This allows us to isolate the contribution of the post-processing framework from detection errors. Initial tracklets are generated using Deep-EIoU \cite{deep-eiou}. For each frame, Deep-EIoU matches the provided ground truth detections with existing tracks using a combination of Expanded IoU and appearance similarity. ReID embeddings are extracted using OSNet x1.0 \cite{osnet} trained on the SportsMOT dataset, which matches the configuration used by Sun et al. \cite{gta} to enable reproducibility. These embeddings are stored per tracklet and used during the tracklet refinement stages. The trackers hyperparameters are tuned on the validation split of the SoccerNet Player-Tracking dataset using grid search over the detection threshold and track buffer length, while the other parameters use the default values reported in the original paper.



\section{Attribute Prediction}\label{attribute_prediction}
The tracklet refinement module relies on predicted domain-specific attributes. Two main attributes need to be predicted: jersey numbers and team identities. 



\subsection{Jersey number prediction}
Jersey number prediction is based on the pipeline from Koshkina et al. \cite{jersey-prediction}. The pipeline processes each tracklet separately and predicts a jersey number including an entropy value that measures prediction uncertainty. We adopt  the same overall structure, but omit the main subject filtering step from the original pipeline,  as it hindered attribute prediction that degraded the tracklet splitters performance. We also transferred the digit compatibility check to the tracklet splitter instead of the prediction pipeline. The remaining stages are described below and a overview can be seen in Figure \ref{fig:jersey_prediction}.

\begin{figure}[h!]
\centering
\includegraphics[width=0.8\textwidth]{figures/jersey_pipeline.png}
\caption{Jersey number prediction pipeline \cite{jersey-prediction}}
\label{fig:jersey_prediction}
\end{figure}

% \subsubsection{Main subject filtering}
% In football broadcast footage, players are often occluded by other players. As a result, some crops may contain appearance information from more than one identity. This could create mistakes in predicting attributes. To address this, Koshkina et. al \cite{jersey-prediction} proposed a main subject filter based on Centroid-ReID \cite{centroid-reid} embeddings. For each tracklet, all crops are used to extract appearance embeddings and the centroid of the embedding distribution is calculated. Crops whose euclidean distance from the centroid exceeds the mean distance by more than a fix threshold are treated as outliers and discarded.

% While Koshkina et al. describe the filter to exclude crops whenever their euclidean distance exceeds the mean distance by $N$ times the standard deviation, their released implementation uses a predefined threshold instead. In our experiments, using the standard deviation variant produced worse filtering results, while the fixed threshold formulation performed well. 



\subsubsection{Legibility classifier}
Not every crop in a tracklet contains a readable jersey number. Since jersey numbers in football are often located on the back of the jersey, a crop is only useful when the player is facing away from the camera. Even then, the number may be partially occluded by another player or the number may be unreadable due to motion blur. Passing such crops to the Scene Text Recognition (STR) model would introduce false and misleading prediction that corrupt the tracklet-level jersey number predictions. To prevent this, Koshkina et al. \cite{jersey-prediction} proposed a legibility classifier: a ResNet34 binary classifier trained to predict whether a torso crop contains a readable jersey number. Crops classified as non-legible are discarded and not used for text recognition.

\subsubsection{Pose detector}
A pose detector is applied to each remaining crop to extract the torso region containing the jersey number. ViTPose \cite{vitpose} estimates COCO-format body keypoints, where the torso region is defined by the left shoulder, right shoulder, left hip and right hip keypoints. A small padding is added to the top and sides to ensure that the full jersey number area is included in the crop.

If any of the four keypoints obtain a confidence score below $0.4$, the pose estimate is considered unreliable and the crop is discarded. Low keypoint confidence typically indicates that the torso is not clearly visible, for example due to occlusion or a poor detection. Koshkina et al. fall back in this case to a simple heuristic that crops a fixed vertical region between $20\%$ and $80\%$ of the bounding box, we instead drop these crops entirely, and not use them for jersey number prediction.


\subsubsection{Scene text recognition}
The remaining torso crops are passed to PARSeq \cite{parseq}, a state-of-the-art scene text recognition model. Since football jersey numbers consist solely of the digits 0-9, Koshkina et. al propose to slice the logits to the first three sequence positions: tens digit, units digit and the End-of-String (EOS) token and to the first eleven vocabulary tokens $E0123456789$, where $E$ denotes the EOS token. A softmax is then applied over this vocabulary, preventing the model from assigning probability scores to irrelevant characters.

For each crop, the confidence is computed as the product of the probabilities of each digit, and the prediction uncertainty is measured using the Shannon entropy of the probability distributions of each digit:

\[
H = - \sum_{k} p_k\log p_k
\]
where $p_k$ is the predicted probability for digit class $k$. 


\subsection{Team classification}
Team identity prediction is formulated as an unsupervised clustering problem within each sequence. In a football match the relevant identities are the two teams, while the numeric team labels themselves are arbitrary. Therefore, no supervised team classifier is trained. Instead, all valid torso crops from all tracklets in the sequence are embedded jointly and clustered into two groups. The resulting cluster labels are only interpreted within the same sequence: cluster 0 in one sequence does not need to correspond to cluster 0 in another sequence.

The team classifier reuses the torso crops extracted by the jersey pipeline. However, unlike jersey number prediction, these crops are taken before the legibility classifier is applied. This is important because team colours can often be recognized from the front or side of the player, whereas jersey numbers are usually only readable from the back. Frames for which no reliable torso crop can be extracted by the pose detector are left without a team prediction.

For each remaining torso crop, a SigLIP vision model is used to extract a 768-dimensional visual embedding. Let $\mathbf{e}_i \in \mathbb{R}^{768}$ denote the SigLIP embedding for crop $i$. Since the team separation is mostly determined by global appearance and colour, the embeddings are first projected to a low-dimensional space using UMAP with cosine distance. In the implementation, UMAP maps the embeddings to three dimensions using $30$ neighbours, a minimum distance of $0$, and a fixed random seed for reproducibility. KMeans is then applied to the projected embeddings with $k=2$, producing one cluster for each team.

The three-dimensional UMAP projection is used as a compromise between visualization-style compression and preserving enough structure for clustering. A two-dimensional projection is convenient for plotting, but it can merge visually distinct modes that are still separable in the original SigLIP space, for example due to lighting, pose, or goalkeeper appearances. Higher-dimensional projections preserve more information, but also retain more nuisance variation and make the following KMeans step closer to clustering the raw high-dimensional embeddings. Three dimensions therefore provide one extra degree of freedom compared to a 2D visualization while still enforcing a compact low-dimensional representation in which colour-based team structure dominates.

KMeans is chosen because the match contains exactly two player teams and because the UMAP projection is intended to make these teams form compact appearance clusters. This does impose a modelling assumption: KMeans assigns each crop to the nearest centroid and therefore works best when the clusters are approximately compact, convex, and similarly scaled. It does not model elongated or unequal-variance clusters explicitly. A Gaussian Mixture Model (GMM) is a natural alternative because it can model elliptical clusters and produces soft assignments, but these probabilities are still not guaranteed to be calibrated after UMAP and can be sensitive to density artefacts in the projection. We therefore use KMeans as a simple and reproducible clustering rule, and handle uncertainty through the centroid-distance confidence score below.

Let $\mathbf{x}_i$ denote the UMAP-reduced embedding of crop $i$, and let $c_i \in \{0,1\}$ be the KMeans cluster assigned to this crop. Because KMeans does not provide calibrated probabilities, a distance-based confidence score is computed from the relative distance to both cluster centroids. Let $\boldsymbol{\mu}_{c_i}$ be the centroid of the assigned cluster and $\boldsymbol{\mu}_{1-c_i}$ the centroid of the other cluster. The confidence is defined as:

\begin{equation}\label{eq:team_conf}
    \operatorname{conf}_T(\mathbf{x}_i)=
    \frac{d(\mathbf{x}_i,\boldsymbol{\mu}_{1-c_i})}
    {d(\mathbf{x}_i,\boldsymbol{\mu}_{c_i}) + d(\mathbf{x}_i,\boldsymbol{\mu}_{1-c_i}) + \epsilon}
\end{equation}

where $d(\mathbf{x}, \boldsymbol{\mu})$ is the Euclidean distance between the projected embedding and a cluster centroid, and $\epsilon$ is a small constant for numerical stability. Since KMeans assigns each point to the nearest centroid, this score lies close to $0.5$ when the crop is near the decision boundary and approaches $1$ when the crop is much closer to its assigned centroid than to the other centroid. The score is used as a reliability measure rather than as a calibrated probability.

Finally, the crop-level predictions are mapped back to the original tracklet frames. Frames with a valid torso crop receive the predicted team label, the distance-based confidence, and the original SigLIP embedding. Frames without a valid crop receive no team prediction and are ignored by the downstream team splitter and tracklet-level team aggregation.



\section{Tracklet Splitter}\label{tracklet_splitter}
The tracklet splitter is designed to split tracklets that contain more than one identity into smaller fragments. An identity switch occurs when the Deep-EIoU tracker incorrectly assigns detections from two different players to the same tracklet, often during occlusion. The framework uses two domain-specific signals to detect these switches: jersey number predictions, and team identity predictions. 

\

The tracklet splitter uses the following core logic: for each tracklet, scan the per-frame predicted attributes for a change in identity. Whenever a change is detected, the framework uses a lookahead window to confirm whether the change is persistent or noise. A detected change is considered to be persistent if the new value appears a sufficient number of times within the window and is the dominant prediction.  This persistence change requirement is critical because attribute predictions are noisy, for example jersey numbers are often unreadable due to (partial) occlusion or motion blur, and team predictions can flip due to occlusion from a player from the opposing team. Splitting on every detected change would produce too many fragments. In this section, the precise working of the tracklet splitter is explained. An overview of the tracklet splitter is shown in Figure \ref{fig:tracklet_splitter_overview}.


\begin{figure}[H]
\centering
\includegraphics[width=0.85\textwidth]{figures/tracklet_splitter_overall_pipeline.drawio.png}
\caption{Overview of the tracklet splitter}
\label{fig:tracklet_splitter_overview}
\end{figure}

\paragraph{Reference Establishment}
Before the scan for identity switches starts, an initial "ground truth" identity must be established independently for each signal. This reference identity represents the current known identity of the player in the tracklet and serves as the baseline to which new predictions are compared. The system scans from the first frame forward to find the earliest reliable prediction that is also persistent. 

\

Let $\mathcal{J} = \{j_1, j_2, \dots ,j_n\}$ denote the per-frame jersey predictions and $\mathcal{H} = \{h_1, h_2, \dots, h_n\}$ be the corresponding Shannon entropy, where $n$ is the tracklet length. A jersey number prediction is often made on a blurry or occluded crop which makes predictions often unreliable. Therefore, a prediction $j_i$ is considered reliable if its entropy is sufficiently low:

\begin{equation} \label{valid_j}
    \text{valid}_J(i) \iff j_i \neq \text{null} \land h_i \leq \tau_J
\end{equation}

The reference jersey number $j_{ref}$ is the first reliable prediction that passes the persistence check, ensuring that the reference is a valid prediction and not noise. Similarly, let $\mathcal{T}=\{t_1,t_2,\dots,t_n\}$ denote the per-frame team predictions and $\text{conf}_T=\{\text{conf}_1,\text{conf}_2,\dots,\text{conf}_n\}$ be the corresponding confidence. 

\begin{equation} \label{valid_t}
    \text{valid}_T(i) \iff t_i \neq \text{null} \land \text{conf}_T(i) \geq \tau_T
\end{equation}

The reference team identity $t_{ref}$ is established using the same procedure as the jersey number reference and is described above.

\paragraph{Persistence Check}
The persistence check separates real identity switches from noise. Instead of requiring the new identity to dominate the full lookahead window, the algorithm only requires the new identity to dominate in the sequence of valid predictions immediately following the switch point. For example, within a lookahead window of $W_j=100$ frames, a new jersey might dominate the first 30 frames of the window before changing back to the original player for the last 70 frames. A global dominance check would see that the original player dominates the window and therefore would not split the tracklet into fragments. By stopping prediction collection as soon as a total of $P_j=20$ reliable jersey number predictions are found, the check only evaluates these predictions and successfully triggers a split.

\

For jersey numbers, the algorithm collects the first $P_j$ reliable predictions after the candidate switch point at frame $i$, within a maximum lookahead window of $W_j$ frames. Let $f_{1},f_{2},\dots,f_{P_j}$ be the strictly increasing sequence of frames such that $f_1 \geq i$ and $f_{P_j} < i+W_j$, and every prediction $\text{valid}_J$ holds. Then we can define the jersey lookahead window as:

\begin{equation}
    \mathcal{W}_J(i)=(j_{f_1}, j_{f_2}, \dots,j_{f_{P_j}})
\end{equation}

A candidate jersey number $c_j$ is considered persistent starting from frame $i$ if it accounts for at least a fraction $\rho_j$ of these predictions:

\begin{equation}
    \text{persistent}_J(i, c_j) \iff 
    |\mathcal{W}_J(i)| = P_j \land 
    \frac{|\{j \in \mathcal{W}_J(i) : j=c_j\}|}{P_j} \geq \rho_j
\end{equation}

The absolute count requirement $P_j$ serves as protection against sparse data. Without it, a high ratio could be obtained by only a small number of valid predictions. For example, if the algorithm only finds 5 of the same valid predictions within the maximum lookahead window, the ratio would be 100\%. Requiring $|\mathcal{W}_J(i)|=P_j$ ensures that the splitter will not trigger on noisy jersey number predictions. 

\

For team identity, predictions are available for almost every frame with a valid torso crop, since KMeans assigns a team identity to every projected SigLIP embedding. However, not all team predictions are equally reliable. A distance-based confidence $\text{conf}_t(i)$ is computed for each prediction as shown in Equation \ref{eq:team_conf}. The algorithm collects all predictions within a fixed lookahead window of $W_t$ frames after the candidate switch point at frame $i$. Let $K$ denote the number of valid team predictions found within this window, and let $f_{1},f_{2},\dots,f_{K}$ be the strictly increasing sequence of frames such that $f_{1} \geq i$ and $f_{K} < i + W_t$ , and every prediction $\text{valid}_T$ holds. Then we can define the team lookahead window:

\begin{equation}
    \mathcal{W}_T(i)=(t_{f_1},t_{f_2},\dots,t_{f_K})
\end{equation}

A candidate team prediction $c_t$ is considered persistent starting from frame $i$ if it accounts for at least a fraction $\rho_t$ of the predictions in the lookahead window:

\begin{equation}
    \text{persistent}_T(i,c_t) \iff |\mathcal{W}_T(i)| > 0 \land  \frac{|\{t \in \mathcal{W}_T(i):t=c_t\}|}{{|\mathcal{W}_T(i)|}} \geq \rho_t
\end{equation}

Unlike the jersey persistence check, no absolute count requirement is needed. Because valid team predictions exist for almost every frame, the lookahead window contains enough predictions for the ratio to be meaningful on its own.

\paragraph{Jersey Number Compatibility}
Before a candidate jersey number detection can trigger a split, it must first pass a digit compatibility check. Jersey numbers often contain two digits. Occlusion or player pose can cause the PARSeq model to only predict one of the two digits. These partial reads would trigger false splits if they are treated as different identities. Therefore, the jersey number compatibility rule states that if one prediction is a single digit that appears as a substring of the other two digit number, then the two predictions are compatible:

\begin{equation}
    \begin{split}
    \text{compatible}(a,b) \iff &
    (|a| =1 \land |b|=2 \land a \subset b) 
    \lor \\
    & (|a| =2 \land |b|=1 \land b \subset a)
    \end{split}
\end{equation}


where $a$ and $b$ denote the detected jersey numbers and $|a|$ and $|b|$ denote the amount of digits. 


\paragraph{Frame-by-Frame Scan and Split Decision}
Once both references are established, the algorithm performs a single scan over all the frames. At each frame, it independently evaluates the jersey and team predictions against the established references. The jersey number splitter triggers when a valid prediction differs from the established jersey reference, passes the digit compatibility check described below, and satisfies the jersey persistence check:

\begin{equation}
    \text{trigger}_J(i) \iff \text{valid}_J(i)  \land j_i \neq j_{ref} \land \text{persistent}_J(i, j_i) \land \neg \text{compatible}(j_i,j_{ref}) 
\end{equation}

The team splitter triggers when a valid prediction differs from the established team reference and satisfies the team persistence check.

\begin{equation}
    \text{trigger}_T(i) \iff \text{valid}_T(i)  \land t_i \neq t_{ref} \land \text{persistent}_T(i, t_i)
\end{equation}

A split point is created if one of the splitters triggers:

\begin{equation}
    \text{split}(i) \iff \text{trigger}_J(i) \lor \text{trigger}_T(i)
\end{equation}

The tracklet is then split at the frame $i$, the first frame that the new identity was detected, rather then at the end of the lookahead window. This ensures that the split point is as close as possible to where the actual identity switch occurred. After the split point is recorded, the reference identity of the signal that caused the split is updated to reflect the new identity:

\begin{equation}
    j_{ref} = 
    \begin{cases}
        j_i, & \text{if trigger}_J(i) \\
        j_{ref}, & \text{otherwise}
    \end{cases}
    \qquad
    t_{ref} = 
    \begin{cases}
        t_i, & \text{if trigger}_T(i) \\
        t_{ref}, & \text{otherwise}
    \end{cases}
\end{equation}

\paragraph{Short Fragment Filtering}
After the frame-by-frame scan is complete, the recorded split points break the tracklets into smaller fragments. Very short tracklets are unlikely to contain genuine player identities, but are often noisy predictions. Each fragment $\text{frag}_m$ is therefore only retained if it contains at least $L$ number of frames:

\begin{equation}
    \text{keep}({\text{frag}_m}) \iff |\text{frag}_m| \geq L
\end{equation}

where $|\text{frag}_m|$ denotes the number of frames in fragment $\text{frag}_m$. If there are fewer than two tracklets which satisfy this condition, then no splits are performed and the original tracklet is returned.




\subsection{Explored Alternatives}
During development, multiple tracklet splitters were explored and evaluated, such as a bounding box anomaly, trajectory and a Global Tracklet Association based tracklet splitter. None of these splitters are used in the final version of the tracklet refinement module, since they did not improve the performance.



\subsubsection{}{Spatio-temporal ReID Splitter}\label{reid_spatio_temporal_splitter}
Modern trackers still have difficulties with disappearing and reappearing players, and frequently create identity switches when a player leaves the camera view and another player reappears at approximately the same position. This causes a single track ID to contain more than one player. The attribute-based splitter can only detect identity switches when these attribute predictions are available for both identities. If they are not, the identity switch will not be detected. The spatio-temporal ReID splitter however, detects identity switches, relying only on temporal discontinuities and the raw ReID embeddings extracted by the Deep-EIoU tracker. 

\paragraph{Temporal Gap Detection}
Let $\mathcal{F}_A=(f_1,f_2,\dots,f_n)$  be the strictly increasing sequence of frame indices of tracklet $A$, and $\textbf{e}_i$ the corresponding ReID embedding at frame $f_i$. A temporal gap is then defined as a position $k$ in this sequence where the distance to the next frame index exceeds a threshold $\tau_{gap}$:

\begin{equation}
    \text{gap}(k) \iff f_{k+1} - f_k \geq \tau_{gap}
\end{equation}

We only consider gaps of at least $\tau_{gap}=5$ frames, smaller gaps are treated as missed detections of the same player and not as candidate split points.

\paragraph{Appearance Comparison}
For every detected gap,  the splitter compares the appearance of the player in the segment before the gap to the segment directly after the gap. To avoid that a single noisy embedding dominates the comparison, we use the local means over a sampling window of $N$ frames on each side. Let $k$ be the position of the detected gap, we can define the embeddings:

\begin{equation}
    \textbf{e}_k^\text{before} = \frac{1}{|B_k|}\sum_{i \in \mathcal{B}_k}\textbf{e}_i,
    \quad 
    \textbf{e}_k^\text{after} = \frac{1}{|\mathcal{A}_k|}\sum_{i\in \mathcal{A}_k}\textbf{e}_i
\end{equation}

where $\mathcal{B}_k=\{\max(1, k-N+1),\dots, k\}$ and $\mathcal{A}_k=\{k+1,\dots,\min(n,k+N)\}$ are the frame indices on either side of the gap. The appearance dissimilarity across the gap is then measured by the cosine distance between these two means:

\begin{equation}
    d_\text{reid} (k)=1-\frac{\textbf{e}_k^\text{before} \cdot \textbf{e}_k^\text{after}}{||\textbf{e}_k^\text{before}|| \cdot ||\textbf{e}_k^\text{after}||}
\end{equation}

\paragraph{Split Decision}
An identity switch is recorded when the cosine distance exceeds a fixed threshold $\tau_\text{reid}$:

\begin{equation}
    \text{split}(k) \iff \text{gap}(k) \land d_\text{reid} (k) > \tau_\text{reid}
\end{equation}

All positions that satisfy this condition are collected into a set of split points. The tracklet is then cut at each split point, producing the resulting fragments.

\paragraph{Short Fragment Filtering}
Very short fragments produced by the spatio-temporal ReID splitter are unlikely to have a reliable identity. Therefore, each fragment is only retained if it contains at least $L_{reid}$ frames:

\begin{equation}
    \text{keep}(\text{frag}) \iff |\text{frag}| \geq L_{reid}
\end{equation}

If fewer than two fragments of sufficient length remain, the split is discarded and the tracklet remains unchanged.



\paragraph{Bounding Box Anomaly Splitter}\label{bbox_anomaly_splitter}
The first approach detects identity switches by identifying abnormal velocity spikes in bounding box detections. The intuition is that when the tracker switches identity, the bounding box teleports from one player to another, causing a high velocity spike inside the tracklet. Let $\textbf{c}_i$ denote the center of the bounding box of the tracklet at frame $i$.

\begin{equation}
    \textbf{c}_i = \left( (\frac{x_1^i + x_2^i}{2},\ \frac{y_1^i + y_2^i}{2}\right)
\end{equation}

Then the velocity magnitude is calculated as:

\begin{equation}
    v_i=||\mathbf{c}_i - \mathbf{c}_{i-1}||_2
\end{equation}

where $\textbf{c}_i$ is the center of the bounding box at frame $i$. A bounding box anomaly is detected when the current velocity exceeds the local baseline by more than $\sigma$ standard deviations and also exceeds a minimum threshold $v_{min}$:

\begin{equation}
    \text{anomaly}(i) \iff \left( \frac{v_i - \mu_{local}}{s_{local}} > \tau \right) \land (v_i > v_{min}) 
\end{equation}

where $\mu_{local}$ is the mean and $s_{local}$ is the standard deviation of the velocities inside the lookback window. To distinguish actual  identity switches from real fast moving players, the spike must last no more than $d_{max}$ frames and the velocity  must return to values close to the baseline. In practice, this approach produced too many false positives since football players move very fast and in many directions.


\paragraph{Trajectory Splitter}
The second approach uses pairwise trajectory analysis and spatial proximity detection. Unlike the bounding box anomaly splitter described in Section \ref{bbox_anomaly_splitter}, the  trajectory splitter detects identity switches by analyzing whether two nearby tracklets exchange trajectory. Two tracklets are considered to be close when they have significant bounding box overlap or the center of both bounding boxes is below a certain proximity threshold $d_{prox}$:

\begin{equation}
    \text{proximate}(A,B,i) \iff \text{IoU}(\text{bbox}_A^i, \text{bbox}_B^i) > \varepsilon_{\text{iou}} \lor ||\textbf{c}_A^i - \textbf{c}_B^i||_2 < d_{prox}
\end{equation}

At each proximity event, three signals are combined to calculate an evidence score: a velocity swap score measuring whether two tracklets exchange velocity directions, a direction change score measuring the angular change in each trajectory, and a speed change score measuring  the  change in velocity magnitude. The trajectory splitter is triggered when the combined evidence score exceeds a threshold $\tau_\text{score}$:

\begin{equation}
    \text{split}(A,B) \iff \text{score}(A,B) \geq \tau_{\text{score}}
\end{equation}

In practice, this approach produced too many false positives since events such as tackles produce natural velocity exchanges without the occurrence of identity switches.

\paragraph{GTA Splitter}
The third splitter uses a DBSCAN \cite{dbscan} to detect identity switches based on ReID features. This approach is adapted from the tracklet splitter method, described in GTA \cite{gta}. Given embeddings for every frame $\text{\textbf{e}}_1, \text{\textbf{e}}_2,\dots, \text{\textbf{e}}_n$, DBSCAN is applied with the cosine distance, neighborhood radius $\varepsilon$, and minimum samples $s$. An identity switch is detected when more than one cluster is found:

\begin{equation}
    \text{switch detected} \iff |\text{\{clusters\}}|>1
\end{equation}

This approach fails often for identity switches that occur with players from the same team, since the ReID features for these players are very similar. This splitter also over-splits clean tracklets and therefore creates too many false positives.

% \subsection{Learned Tracklet Splitter}

\section{Tracklet Merger}\label{tracklet_merger}
After the tracklet splitting stage, we arrive at the tracklet merger. The tracklet splitter outputs a set of fragmented tracklets that ideally contain a single identity. The role of the tracklet merger is to reconnect fragments belonging to the same player into a single complete trajectory. 

\

The tracklet merger makes use of the predicted jersey numbers and team identities, spatio-temporal information, and appearance features in order to connect tracklets containing the same identities. However, jersey number and team identity predictions can be missing or unreliable, making the merger only able to rely on ReID appearance features, which is considerably harder given the visual similarity between players.

\

Three merger designs are developed and evaluated: a rule-based merger, a gradient boosted trees merger and a transformer-based merger. Before describing each design, we first define how per-frame predictions are aggregated to the tracklet level. Since individual per-frame predictions can be noisy, we apply confidence-based filtering and then analyze whether a tracklet contains a reliable jersey number and team identity.

\

\subsection{Jersey Number Aggregation}
Not all jersey number predictions are equally reliable. For each tracklet $A$ with frames $\mathcal{F_A}$, the set of reliable jersey number predictions is defined as all frames passing the validity check from Equation \ref{valid_j}:

\begin{equation}
    \mathcal{V}_J(A) = \{j_i : i \in \mathcal{F}_A \land \text{valid}_J(i)\}
\end{equation}

A tracklet has a reliable jersey number identity when the tracklet contains at least $N_J^{min}$ predictions. When this conditions holds, the tracklet-level jersey number identity is the most frequently occurring prediction within its tracklet, also known as the mode.

\begin{equation}
    \text{Identity}_J(A) = \text{mode}(\mathcal{V}_J(A)) \iff |\mathcal{V}_J(A)| \geq N_J^{min}
\end{equation}

\subsection{Team Identity Aggregation}
Similarly, the set of reliable team identity predictions is defined as all frames passing the validity check from Equation \ref{valid_t}:

\begin{equation}
    \mathcal{V}_T(A) = \{t_i : i \in \mathcal{F}_A \land \text{valid}_T(i)\}
\end{equation}

Tracklet-level team identity does not require a minimum count of predictions, since these predictions are available on almost every frame. Instead, team predictions need to be consistent. Therefore, we define the team consistency requirement to measure how dominant the predominant team is among the reliable predictions:

\begin{equation}
    \gamma_T(A)=\frac{|\{v \in \mathcal{V}_T(A):v=t_A^*\}|}{|\mathcal{V}_T(A)|}
\end{equation}

A tracklet has a reliable team identity when the team predictions are consistent and thus confident enough to exceed a threshold $\theta_T$:

\begin{equation}
    \text{Identity}_T(A) = \text{mode}(\mathcal{V}_T(A)) \iff \mathcal{V}_T(A) \neq \emptyset \land \gamma_T(A) \geq \theta_T
\end{equation}

Team identity uses two-stage filtering: first by per-frame confidence, then by tracklet-level consistency. This ensures that team identity is only trusted and used for merging when the team classifier is confident on the predicted identities and those confident predictions agree with each other.

\subsection{Rule-based Decision Merger}
The rule-based decision merger evaluates each pair of tracklets through a hierarchical decision tree that prioritizes the most discriminative signals. Each tracklet pair is classified into one of the three branches: merge, block, or ReID. This structure ensures that confident predicted identities are always prioritized over visual similarity, while tracklets that lack the identity information can still be merged using ReID appearance features.




\paragraph{Pairwise Decision Tree}
The pairwise decision tree works as follows, for a candidate pair $(A,B)$, the merger applies the following rules in order:

\begin{enumerate}
    \item \textbf{Temporal overlap}. If tracklets $A$ and $B$ contain detections in overlapping frames $(\mathcal{F}_A \cap \mathcal{F}_B \neq \emptyset)$, they can never be the same player $\rightarrow$ decision: BLOCK
    \item \textbf{Jersey conflict}. If both tracklets $A$ and $B$ have confident but conflicting jersey numbers $\text{Identity}_J(A) \neq \text{Identity}_J(B)$, then the tracklets belong to different players $\rightarrow$ decision: BLOCK
    \item \textbf{Full identity match}. If two tracklets both have the same confident jersey number and confident team identity $(\text{Identity}_J(A) = \text{Identity}_J(B) \land \text{Identity}_T(A) = \text{Identity}_T(B))$, then these tracklets belong to the same player $\rightarrow$ decision: MERGE
    \item \textbf{Jersey match, team conflict}. If two tracklets have the same confident jersey number but their team identity  disagree $(\text{Identity}_J(A) = \text{Identity}_J(B) \land \text{Identity}_T(A) \neq \text{Identity}_T(B))$, the tracklets belong to players with the same jersey number but from opposing teams $\rightarrow$ decision: BLOCK
    \item \textbf{Jersey or team unclear}. If at least one of the tracklets does not have a confident jersey number or team identity, then we can not conclude whether the two tracklets belong to the same player using the attributes alone $\rightarrow$ decision: REID
\end{enumerate}


\begin{figure}[h!]
\centering
\includegraphics[width=1\textwidth]{figures/decision_merger.drawio.png}
\caption{Overview of the pipeline}
\label{fig:decision_merger_overall_pipeline}
\end{figure}




\paragraph{ReID Fallback}
For a tracklet pair classified as REID, the merger computes the average pairwise cosine distance between all ReID embeddings of the two tracklets. The pair is eligible for merging only if the cosine distance is smaller than a predefined value $\theta_{reid}$. For each embedding in tracklet $A$, the cosine similarity is computed against every embedding in tracklet $B$. The resulting $n_A \times n_B$ similarities are then averaged into a single distance score, making the measure more robust to noisy or occluded frames. The pairwise cosine ReID distance is calculated as follows:

\begin{equation}\label{eq:cos_reid_dist}
    d_{reid}(A,B)=1 - \frac{1}{n_A \cdot n_B} \sum_{i=1}^{n_A}\sum_{k=1}^{n_B} \frac{\textbf{e}_i^A \cdot \textbf{e}_k^B} {||\textbf{e}_i^A|| \cdot ||\textbf{e}_k^B||}
\end{equation}

where $\textbf{e}_i^A$ denotes the ReID embedding extracted at the i-th detection of tracklet A, and $n_A$ and $n_B$ are the number of embeddings in each tracklet.




\paragraph{Greedy Merge}
Each iteration, the merger makes greedy choices. It evaluates all pairs, applies the decision tree, and selects the pair with the lowest distance:

\begin{itemize}
    \item Pairs classified as MERGE receive a distance $d=0$, ensuring that they are always merged first.
    \item Pairs classified as REID receive a distance corresponding to their computed cosine ReID distance $d_{reid}$
    \item Pairs classified as BLOCK are excluded entirely and  are never merged.
\end{itemize}

The selected pair is then merged by appending all data from one tracklet into the other, and the appended tracklet is removed from the set. The loop continues until no eligible pairs remain. This greedy strategy processes all pairs classified as MERGE first, since their distance $d=0$ is the minimum, followed by merging the REID classified pairs where the pairs with the lowest $d_{reid}$ are merged first. 



\subsection{XGBoost Merger}
The rule-based decision merger uses predetermined thresholds and rules to determine if two tracklets should be merged. This merger performs well when confident attribute predictions are available, but it falls back to a simple ReID threshold when they are not. This approach also fails to capture complex interactions between features. For example, a moderate  ReID similarity combined with matching jersey numbers and a small temporal gap may strongly suggest that it is the same player, but the rule-based merger assesses each signal independently and is unable to learn such patterns. We propose an XGBoost merger that addresses this limitation by replacing the hand-crafted rules and thresholds with a supervised gradient boosted trees model that learns to predict the probability that two tracklets belong to the same player.

\begin{figure}[h!]
\centering
\includegraphics[width=0.8\textwidth]{figures/xgboost_example.drawio.png}
\caption{XGBoost example}
\label{fig:gta_splitter}
\end{figure}

\begin{figure}[H]
\centering
\includegraphics[width=0.8\textwidth]{figures/xgboost.drawio.png}
\caption{Overview of the XGBoost merger}
\label{fig:xgboost_overall_pipeline}
\end{figure}


\paragraph{Training Data Generation}
The training data is generated from the tracklets produced by the tracklet splitter, including their predicted attributes. For each sequence, the following procedure is applied:

\begin{enumerate}
    \item \textbf{Tracklet aggregation}. Each tracklet is collapsed into a fixed-size feature vector summarizing its information (described in the next paragraph).
    \item \textbf{Candidate pairs}. All tracklet pairs that do not overlap in time are considered merge candidates.
    \item \textbf{Pairwise features}. For each candidate pair, additional features are computed that capture the relationship between the two tracklets.
    \item \textbf{Labeling}. Each pair is labeled using the ground truth track identities: $\text{label}=1$ if both tracklets share the same ground truth identity, and $\text{label}=0$ if the pair contains two different players.
\end{enumerate}

This produces a CSV file for every data split where each row represents a candidate tracklet pair, containing the label that indicates whether the pair should be merged and all the tracklets features. An example can be found in Appendix \ref{appendix:...}



\paragraph{Class Imbalance}
Each sequence contains many different tracklets, but each player produces only a few fragments. This causes the data to contain much more tracklets with a different identity rather than the same identity. To address this, negative pairs are down-sampled by a fixed ratio. We also utilize XGBoost's \texttt{scale\_pos\_weight} to further penalize missed merges by setting it to $n_{neg}/n_{pos}$.

\paragraph{Tracklet-level Feature Aggregation}
Each tracklet is collapsed into a fixed-size feature vector summarizing its appearance, identity attributes, spatial position, and temporal information. The aggregation produces four features:

\begin{itemize}
    \item \textbf{Jersey number features}. The tracklet-level jersey identity is defined as the most frequently predicted valid jersey number across all frames. The mean Shannon entropy of the jersey number predictions measures the prediction uncertainty, and the jersey coverage measures the fraction of frames that contain a valid jersey number prediction.

    \item \textbf{Team identity features}. The tracklet-level team identity is defined as the most frequently predicted valid team identity across all frames. The mean confidence of the team identity predictions measures the prediction confidence, and the team coverage measures the frames that contain a valid team prediction.

    \item \textbf{Temporal features}. Temporal features include the start frame, end frame, and total duration of the tracklet.

    \item \textbf{Spatial features}. Spatial features include the bounding box coordinates at the first and last frame and the mean bounding box height across all frames. The coordinates capture where the player was seen before and after appearing in the frame. The mean bounding box height helps to identify whether a player appears close to or far from the camera. 
\end{itemize}


\paragraph{Pairwise Feature Computation}
While tracklet-level features describe each tracklet independently, the strongest signals for merging come from pairwise comparisons. For each candidate pair $(T_A,T_B)$, 12 pairwise features are computed that capture temporal, spatial, appearance, and attribute relationships.

\begin{itemize}
    \item \textbf{Temporal gap}. The temporal gap defines the number of frames between when the first fragment ends and the second fragment begins. 

    \item \textbf{Spatial distance}. The Euclidean distance between the last position of the first fragment and the first position of the second tracklet measures the distance between the two fragments: $d_{spatial}(T_A,T_B)=\sqrt{(\text{start}_x^\text{entry} - \text{end}_x^\text{exit})^2 + (\text{start}_y^\text{entry}-\text{end}_y^\text{exit})^2}$. We also include the directional displacements $\Delta x$ and $\Delta y$ are included for the model to learn movement patterns, such as the fact that when a player who exits the left side of the camera view will likely reappear from the same side.

    \item \textbf{Bounding box height ratio}. The ratio of mean bounding box heights between the two fragments can indicate whether two tracklets appear at a similar depth on the pitch.

    \item \textbf{Appearance similarity}. The appearance of two fragments is compared using a cosine similarity score. The cosine similarity between two tracklets is computed using the same all-pairs approach as in the rule-based merger seen in Equation \ref{eq:cos_reid_dist}. This feature captures the general player appearance similarity.

    \item \textbf{Jersey number agreement}. Jersey number relationships are encoded using three features: \texttt{jersey\_match}, \texttt{jersey\_conflict}, and \texttt{jersey\_both\_confident}. The jersey both confident feature tells the model whether to trust the match or conflict signal. Both the jersey match and jersey conflict features are needed to encode when one or both tracklets have no jersey predictions.

    \item \textbf{Team identity agreement}. Team identity relationships are encoded using the three features: \texttt{team\_match}, \texttt{team\_conflict}, and \texttt{team\_both\_consistent}.
\end{itemize}


\paragraph{Model Architecture and Training}
The XGBoost classifier is trained on the constructed training dataset including the tracklet-level and pairwise features. The model is evaluated on the validation split using AUC as the evaluation metric. The model uses 1000 decision trees and employs early stopping after 30 rounds to prevent overfitting. Each tree is trained on a random subset of rows and features, which helps to reduce correlation between trees and improves generalization. A minimum loss reduction threshold \texttt{gamma} is required before any split is accepted in order to prevent the model from creating meaningless splits. A minimum child weight is added to ensure that leaves contain enough samples, preventing the model from memorizing noise in the training data. We use positive weight class weight scaling to correct for the class imbalance by increasing the relative penalty for missed merges.

\

Since the XGBoost merger is the best performing variant, we perform systematic hyperparameter tuning to ensure the model is optimally configured. The tuning is performed in two stages: the first stage consists of tuning the XGBoost model hyperparameters by optimizing the validation AUC using a grid search, and in the second stage, the merge threshold $\delta_{\text{merge}}$ is tuned by fixing the best model parameters from stage one and optimizing the HOTA score on the validation split. The selected hyperparameters are shown in Table \ref{tab:xgb_hyperparams}.


\begin{table}[ht]
\centering

\begin{tabular}{l c l}
\hline
\textbf{Parameter} & \textbf{Value} & \textbf{Description} \\
\hline
\texttt{n\_estimators}        & 1000                          & Maximum boosting rounds \\
\texttt{max\_depth}           & 6                             & Maximum tree depth \\
\texttt{learning\_rate}       & 0.05                          & Shrinkage rate per round \\
\texttt{subsample}            & 0.8                           & Fraction of rows sampled per tree \\
\texttt{colsample\_bytree}    & 0.8                           & Fraction of features sampled per tree \\
\texttt{min\_child\_weight}   & 5                             & Minimum samples in a leaf node \\
\texttt{gamma}                & 1.0                           & Minimum loss reduction for a split \\
\texttt{scale\_pos\_weight}   & $n_{\text{neg}} / n_{\text{pos}}$ & Class imbalance correction \\
\texttt{eval\_metric}         & AUC                           & Monitoring metric \\
\texttt{early\_stopping\_rounds} & 30                         & Patience for early stopping \\
\hline
\end{tabular}
\caption{XGBoost training hyperparameters.}
\label{tab:xgb_hyperparams}
\end{table}


\paragraph{Inference: from Probabilities to Merged Tracklets}
During inference, we translate both rule-based constraints and the XGBoost model's predicted probabilities into a pairwise distance matrix. We then use this matrix to perform hierarchical clustering \cite{hierarchical_clustering}. This procedure consists of four steps:

\begin{enumerate}
    \item \textbf{Hard constraint filtering}. The first step is to filter each tracklet pair using three hard constraints which are also used in the rule-based decision merger. If tracklets $T_A$ and $T_B$ have overlapping frames, both have confident but conflicting jersey numbers, or both have consistent but conflicting team identities, the pair will never be merged. 
    
    \item \textbf{XGBoost scoring}. For all remaining tracklet pairs, the XGBoost model predicts a merge probability $p(T_A,T_B) \in [0,1]$. Since hierarchical clustering merges clusters with the smallest distance, we need to transform the probabilities into distances $d(T_A,T_B) = 1-p(T_A,T_B)$. A high merge probability results in a low distance, causing tracklets to be close together in the distance space, whereas a low probability results in a high distance and keeps the tracklets apart. Together, steps 1 and 2 create the pairwise distance matrix where the pairs blocked by the hard constraints described in step 1 receive a distance of 2, and the remaining pairs receive the distance calculated by the XGBoost model.
    
    \item \textbf{Hierarchical agglomerative clustering}. The distance matrix is used to perform hierarchical agglomerative clustering with average linkage. Average linkage computes the distance between two clusters as the mean of all pairwise distances between their members. Average linkage is chosen as it is the most robust to noisy pairwise distances, which is important since ReID similarities and attribute predictions can be unreliable. The clustering is stopped when the distance between any two remaining clusters exceeds the threshold $\theta_\text{merge}$, after which each resulting cluster contains the fragments that should be merged into a single final tracklet. 
    
    \item \textbf{Merging}. After clustering, each cluster contains a group of tracklets that need to be merged into a single tracklet. The tracklets within each cluster are first sorted chronologically by their start frame, and then merged sequentially by concatenating all its features. 
\end{enumerate}

\subsection{Transformer-based Merger}
The XGBoost merger successfully learns complex interactions between features, but it relies on aggregated tracklet-level statistics. By compressing an entire tracklet into a fixed-size set of averages and modes, we lose the temporal structure of the data: the order of frames, how appearance changes across the tracklet, and which frames are located at the boundary where the player disappears and reappears. These are the exact signals that humans use to recognize a player across multiple video clips. A transformer-based merger addresses this limitation by processing per-frame features directly, using self-attention to identify which frames are informative and how two frame sequences relate to each other.

\

Three transformer architectures are developed and evaluated, each representing a different approach to comparing two variable-length tracklet sequences. The \textbf{Siamese CLS} encoder processes each tracklet independently through a shared transformer encoder and compares the resulting fixed-size summaries. The \textbf{Cross-Attention} model allows frames from one tracklet to directly attend to frames from the other, enabling frame-level correspondence matching. The \textbf{Hybrid} model combines both approaches: it first builds intra-tracklet context using self-attention, then performs cross-attention comparison on the enriched representations. All three architectures share the same per-frame token construction, positional encoding, and training procedure, and differ only in how they process and compare the two frame sequences.


\paragraph{Training Data Generation}
Training data is generated using the same labeling procedure as the XGBoost connector: all tracklet pairs that do not overlap in time are considered merge candidates, and each pair is labeled based on whether both tracklets share the same ground truth identity. The same 12 pairwise features described in Section \ref{} are computed for every pair. Unlike the XGBoost connector, the training data does not store per-frame features in every row. Since each tracklet appears in multiple candidate pairs, embedding its full frame sequence in every row would cause massive data duplication. Instead, each training sample stores only the sequence name, two tracklet identifiers, the pairwise feature vector, and the ground truth label. The per-frame features are loaded during training from a pre-computed cache stored on disk.

\

The natural distribution of candidate pairs is heavily imbalanced: each sequence contains many different players but each player produces only a few fragments, resulting in far more negative pairs than positive pairs. Negative pairs are downsampled to a fixed 3:1 negative-to-positive ratio per sequence, consistent with the XGBoost training configuration. Tracklets with a ground truth purity below 0.6 are excluded from pair generation to prevent noisy labels from corrupting the training signal.

\paragraph{Synthetic Positive Augmentation}\label{par:synth_augmentation}
Even after downsampling negatives, the absolute number of real positive pairs remains small: the 57 training sequences produce approximately 1,100 real positive pairs. To increase both the total amount of training data and the diversity of positive examples, we augment the dataset with synthetic positive pairs. Synthetic positives are generated by splitting existing tracklets into two fragments and treating each pair of fragments as a positive example, since both originate from the same player.

\

The split points are not chosen randomly. Instead, they are biased toward frames where the ReID embedding changes the most within the tracklet. For each eligible tracklet (those passing the purity threshold and containing at least 20 frames), we compute the per-frame cosine distance between consecutive ReID embeddings. The split point is then sampled with probability proportional to this drift signal, causing synthetic pairs to be created at positions where appearance changes naturally, such as when a player turns or becomes partially occluded. This produces more realistic positive pairs compared to random splitting, since the two fragments will have visibly different appearances despite belonging to the same player. Each eligible tracklet produces up to two synthetic pairs, and the number of synthetic positives is capped at twice the number of real positives. The final negative-to-positive ratio 
